% Part: sets-functions-relations
% Chapter: functions
% Section: basics
% Hindi translation (hi-Deva-IN) of Open Logic commit 9620cc73f9c8e0ad003c514a5d3748f29611c4c0.

\documentclass[../../../include/open-logic-section]{subfiles}

\begin{document}

\olfileid[hi]{sfr}{fun}{bas}
\olsection{फलनों की मूल बातें}

\begin{explain}
\emph{फलन} वह प्रतिचित्रण है जो किसी दिए गए समुच्चय के प्रत्येक !!{element}
को किसी दिए गए (अन्य) समुच्चय के एक निश्चित !!{element} पर भेजता है। उदाहरण के
लिए,~$1$ जोड़ने की संक्रिया एक फलन परिभाषित करती है: प्रत्येक संख्या~$n$ को
एक अद्वितीय संख्या~$n+1$ पर प्रतिचित्रित किया जाता है।

अधिक सामान्य रूप में, फलन युग्म, त्रिक आदि को निवेश के रूप में ले सकते हैं और
किसी प्रकार का निर्गम लौटा सकते हैं। आधारभूत अंकगणित से बहुत-से फलन हमें
परिचित हैं। उदाहरण के लिए, योग और गुणन फलन हैं। वे दो संख्याएँ लेते हैं और
तीसरी संख्या लौटाते हैं।

इस गणितीय, अमूर्त अर्थ में फलन एक \emph{काली पेटी} है: महत्त्व केवल इसका है
कि कौन-सा निर्गम किस निवेश के साथ युग्मित है, न कि निर्गम की गणना करने की विधि
का।
\end{explain}

\begin{defn}[फलन]
\emph{फलन} $f \colon A \to B$,~$A$ के प्रत्येक !!{element} का~$B$ के किसी
!!{element} पर प्रतिचित्रण है।

हम $A$ को~$f$ का \emph{प्रांत} और $B$ को~$f$ का \emph{सहप्रांत} कहते हैं।
समुच्चय~$A$ के !!{element}ों को~$f$ के निवेश या \emph{फलन-कोणांक} कहते हैं,
और~$B$ के जिस !!{element} को फलन-कोणांक~$x$ के साथ~$f$ द्वारा युग्मित किया
गया है, उसे~$f$ का फलन-कोणांक~$x$ पर \emph{मान} कहते हैं, जिसे~$f(x)$ लिखते
हैं।

\emph{परिसर} $\ran{f}$,~$f$ के सहप्रांत का वह उपसमुच्चय है जिसमें किसी
फलन-कोणांक के लिए~$f$ के मान आते हैं; $\ran{f} =
\Setabs{f(x)}{x \in A}$।
\end{defn}

\olref{fig:function} का आरेख फलनों के बारे में सोचने में सहायक हो सकता है।
बाईं ओर का दीर्घवृत्त फलन के \emph{प्रांत} को निरूपित करता है; दाईं ओर का
दीर्घवृत्त फलन के \emph{सहप्रांत} को निरूपित करता है; और एक तीर प्रांत के
किसी \emph{फलन-कोणांक} से सहप्रांत के संगत \emph{मान} की ओर जाता है।

\begin{figure}[H]
  \olasset{assets/diagrams/function.tikz}
  \caption{फलन एक समुच्चय के प्रत्येक !!{element} का दूसरे समुच्चय के
    !!a{element} पर प्रतिचित्रण है। एक तीर प्रांत के किसी फलन-कोणांक से
    सहप्रांत के संगत मान की ओर जाता है।}
  \ollabel{fig:function}
\end{figure}

\begin{ex}
गुणन प्राकृत संख्याओं के युग्मों को निवेश के रूप में लेता है और उन्हें
प्राकृत संख्याओं पर निर्गम के रूप में प्रतिचित्रित करता है; अतः यह
$\Nat \times \Nat$ (प्रांत) से $\Nat$ (सहप्रांत) तक जाता है। इसका परिसर भी
$\Nat$ ही निकलता है, क्योंकि प्रत्येक $n \in \Nat$,~$n \times 1$ है।
\end{ex}

\begin{ex}
गुणन एक फलन है, क्योंकि वह प्रत्येक निवेश---प्राकृत संख्याओं के प्रत्येक
युग्म---को एक ही निर्गम के साथ युग्मित करता है: $\times \colon \Nat^2 \to
\Nat$। इसके विपरीत, किसी प्राकृत संख्या~$n$ को ऐसी वास्तविक संख्या~$x$ पर
प्रतिचित्रित करना जिसके लिए $x^2 = n$ हो, फलन नहीं है, क्योंकि प्रत्येक धन
पूर्णांक~$n$ के दो वर्गमूल होते हैं: $\sqrt{n}$ और~$-\sqrt{n}$। केवल धनात्मक
वर्गमूल लौटाकर हम इसे फलन बना सकते हैं: $\sqrt{\phantom{X}} \colon
\Nat \to \Real$।
\end{ex}

\begin{ex}
जो संबंध किसी कक्षा के प्रत्येक विद्यार्थी को उसके अंतिम श्रेणी-अंक के साथ
युग्मित करता है, वह फलन है---एक ही कक्षा में कोई विद्यार्थी दो अलग-अलग अंतिम
श्रेणी-अंक नहीं पा सकता। जो संबंध किसी कक्षा के प्रत्येक विद्यार्थी को उसके
माता-पिता के साथ युग्मित करता है, वह फलन नहीं है: विद्यार्थियों के माता-पिता
शून्य, दो अथवा दो से अधिक हो सकते हैं।
\end{ex}

\begin{explain}
हर संभावित फलन-कोणांक के लिए फलन का मान क्या है, यह किसी सटीक ढंग से बताकर हम
फलन परिभाषित कर सकते हैं। ऐसा करने के अलग-अलग तरीके हैं: कोई सूत्र देना, मान
की गणना करने की विधि बताना, अथवा प्रत्येक फलन-कोणांक के लिए मानों की सूची
देना। फलनों को जैसे भी परिभाषित करें, हमें यह सुनिश्चित करना चाहिए कि हर
फलन-कोणांक के लिए एक, और केवल एक, मान निर्दिष्ट हो।
\end{explain}


\begin{ex}
मान लीजिए कि $f \colon \Nat \to \Nat$ को इस प्रकार परिभाषित किया गया है कि
$f(x) = x+1$। यह परिभाषा~$f$ को ऐसे फलन के रूप में निर्दिष्ट करती है जो
प्राकृत संख्याओं को निवेश के रूप में लेता है और प्राकृत संख्याएँ निर्गम करता
है। यह हमें बताती है कि प्राकृत संख्या~$x$ दिए जाने पर~$f$ उसका परवर्ती~$x+1$
निर्गम करेगा। इस स्थिति में सहप्रांत $\Nat$,~$f$ का परिसर नहीं है, क्योंकि
प्राकृत संख्या~$0$ किसी भी प्राकृत संख्या की परवर्ती नहीं है। फलन~$f$ का
परिसर सभी धन पूर्णांकों का समुच्चय $\Int^{+}$ है।
\end{ex}

\begin{ex}\ollabel{examplefunext}
मान लीजिए कि $g \colon \Nat \to \Nat$ को इस प्रकार परिभाषित किया गया है कि
$g(x) = x+2-1$। यह हमें बताता है कि~$g$ ऐसा फलन है जो प्राकृत संख्याओं को
निवेश के रूप में लेता है और प्राकृत संख्याएँ निर्गम करता है। प्राकृत
संख्या~$n$ दिए जाने पर~$g$,~$x$ के परवर्ती के परवर्ती का पूर्ववर्ती, अर्थात
$x+1$, निर्गम करेगा।
\end{ex}

\begin{explain}
हमने अभी अलग-अलग \emph{परिभाषाओं} वाले दो फलनों $f$ और $g$ पर विचार किया।
फिर भी ये \emph{एक ही फलन} हैं। आखिरकार, किसी भी प्राकृत संख्या~$n$ के लिए
$f(n) = n+1 = n+2-1 =
g(n)$ है। दूसरे शब्दों में,~$f$ और~$g$ की हमारी
परिभाषाएँ अलग-अलग समीकरणों द्वारा एक ही प्रतिचित्रण निर्दिष्ट करती हैं। अतः
हम फलनों के लिए विस्तारिता के एक सिद्धांत पर निहित रूप से निर्भर हैं,
\[
  \text{यदि }\forall x\, f(x) = g(x)\text{, तो }f = g
\]
बशर्ते कि~$f$ और~$g$ का प्रांत और सहप्रांत समान हो।
\end{explain}

\begin{ex}
हम फलनों को खंडशः भी परिभाषित कर सकते हैं। उदाहरण के लिए, हम
$h \colon \Nat \to \Nat$ को इस प्रकार परिभाषित कर सकते हैं:
\[
h(x) =
\begin{cases}
  \frac{x}{2} & \text{यदि $x$ सम है} \\
  \frac{x+1}{2} & \text{यदि $x$ विषम है।}
\end{cases}
\]
क्योंकि प्रत्येक प्राकृत संख्या या तो सम है या विषम, इस फलन का निर्गम सदैव
प्राकृत संख्या होगा। केवल यह याद रखें कि यदि आप किसी फलन को खंडशः परिभाषित
करते हैं, तो प्रत्येक संभावित निवेश ठीक एक ही स्थिति में आना चाहिए। कुछ
स्थितियों में इसके लिए यह सिद्ध करना आवश्यक होगा कि स्थितियाँ सर्वसमावेशी
और परस्पर अपवर्जी हैं।
\end{ex}

\end{document}
